% Generated by paper4/manuscript/build.py from paper4/paper/section_introduction.md (90c3911735210f59): edit the Markdown, not this file.
\section{Introduction}
A road operator who installs traffic sensors pays for each one, and covers only part of the network. Two decisions follow: how many sensors to install and where, and what each of them should measure.

For the first decision, the usual hope is that instrumented segments inform the segments without a sensor. Spatio-temporal graph models propagate information along the road network ~\cite{wu2019gwnet}, and kriging methods estimate unobserved locations from observed ones ~\cite{wu2021ignnk}. If propagation works, coverage and placement buy accuracy on the whole network, not only where the sensors are.

We test that hope on one task: classifying congestion in three classes (free flow, slowed, congested) 15–120 minutes ahead, on segments that have no sensor. A segment without a sensor is never seen by the model: it is not in the input, in the training loss, in the normalisation or in the class weights. This is the situation of an installation, and it is stricter than the evaluation of the project report that preceded this work ~\cite{motion2026uc04}, where the segments said to be hidden were training targets. The data are the measured speeds of Madrid, London and Melbourne from the Traffic4cast benchmark ~\cite{neun2023t4c}. The reference is the same temporal network as the graph model, Graph WaveNet ~\cite{wu2019gwnet}, with its graph removed.

For the second decision, we ask what the most common urban sensor, the loop detector, measures of congestion, against speed references independent of its own reconstruction: the M30 ring of Madrid, where the same device measures speed; the Madrid arterials, where loops are matched to probe speeds; and Austin, with radar detectors per lane. The earlier report analysed the same measurements; we regenerate its figures and apply three corrections to its analysis.

Every experiment was pre-registered before any test metric was computed, under the method of an earlier study by the same authors ~\cite{devito2026identity}: signed pre-registrations ~\cite{nosek2018prereg}, frozen records generated by script, intervals from a joint resample of seeds and blocks of the network, and two reading scales.

We address four research questions.

\begin{itemize}
\item \textbf{RQ1 (propagation).} To what extent does propagation from instrumented neighbours improve the forecast on segments without a sensor, as the spacing of the network and the instrumented fraction change?
\item \textbf{RQ2 (signal).} How much forecastable information do the labels of instrumented neighbours carry about a segment without a sensor, once calibrated and once the calendar is known?
\item \textbf{RQ3 (where).} How does the rule that places the sensors change the proximity of uncovered segments to a sensor, and what is that proximity worth?
\item \textbf{RQ4 (what to measure).} To what extent does a loop detector measure the level of congestion rather than its occurrence, and how do the matching to road segments and the aggregation of lanes change what it measures?
\end{itemize}

The answers to RQ1–RQ3 are negative, in a precise sense. The graph model gains nothing over the model without spatial exchange, on compact sub-networks at any spacing from 43 to 397 m and on whole cities at any coverage from 10\% to 60\%. The labels do carry a spatial signal: within about 250 m, a segment without a sensor shares its nearest instrumented neighbour's class more often than chance. Calibrated, that signal comes within 1–3 points of the reference and does not pass it; once the hour of the week is known, it adds nothing. Placement changes proximity by a large factor, and that proximity buys no forecast.

The answer to RQ4 is descriptive, because its reading rules were fixed knowing the in-sample values of the earlier report. A loop detector detects congestion but does not measure its level (rank correlation +0.249 against the speed measured by the same device; area under the curve 0.984), with a threshold that transfers between sensors. Moving the reference some 30 m from the loop accounts for about half of its disagreement with probe speeds on arterials, and aggregating lanes hides episodes on one lane.

\textbf{On segments without a sensor, these data carry no forecastable spatial information beyond the calendar at 15–120 minutes. What an operator buys with a sensor is the measurement at that segment, and the choice of what to measure.}

\subsection{Contributions}
\begin{enumerate}
\item \emph{A decomposition of the spatial signal without models.} We measure the agreement of nearby labels, then its value as a calibrated forecast, then its value net of the calendar, on the same segments, origins and resample as the trained models. The agreement is short, and its forecasting value is that of the shared weekly rhythm; this is why a further learned model has nothing to extract.
\item \emph{An inductive evaluation of propagation to segments without a sensor.} Hidden segments never enter input, loss, normalisation or class weights; node sets are whole blocks of the network, with disjoint sets for training, checkpoint selection and evaluation; a positive control shows what the same pipeline reaches when hidden segments are training targets. In this regime propagation gives no measurable gain.
\item \emph{A geometric account of placement.} For three rules at equal numbers of sensors, we report how many uncovered segments lie within the radius of the signal. Placement decides where the network measures, not how well it forecasts elsewhere.
\item \emph{A measurement audit of the loop detector.} On three datasets with an independent speed reference, with a threshold chosen out of sample, block-bootstrap intervals and a rank correlation for the effect of lanes: a loop detects congested episodes but does not measure their level; matching explains about half of its disagreement with probe data on arterials; lane aggregation hides single-lane episodes.
\item \emph{Reproducible artefacts.} Signed pre-registrations, frozen records with the fingerprints of every input, per-cell predictions, a register of corrections, and one script that regenerates every number.
\end{enumerate}

